% =====================================================================
%  Relatório de execução — Verificador de Contestação (Fase 1)
%
%  Compilar com pdfLaTeX (recomendado). Em Overleaf, basta selecionar
%  pdfLaTeX (padrão). Os blocos de código usam o pacote listings com um
%  mapeamento de acentos (literate) para renderizar corretamente sob
%  pdfLaTeX + inputenc utf8.
% =====================================================================
\documentclass[11pt,a4paper]{article}

\usepackage[utf8]{inputenc}
\usepackage[T1]{fontenc}
\usepackage[brazil]{babel}
\usepackage{lmodern}
\usepackage{textcomp}   % \textdegree, \textordmasculine, etc.
\usepackage{upquote}    % aspas retas em verbatim/listings (upquote=true)
\usepackage[a4paper,margin=2.3cm]{geometry}
\usepackage{xcolor}
\usepackage{listings}
\usepackage{booktabs}
\usepackage{longtable}
\usepackage{array}
\usepackage{enumitem}
\usepackage{parskip}
\usepackage[hidelinks]{hyperref}

% ---------- cores ----------
\definecolor{corFundo}{RGB}{248,248,248}
\definecolor{corMoldura}{RGB}{220,220,220}
\definecolor{corComent}{RGB}{0,128,0}
\definecolor{corChave}{RGB}{0,0,180}
\definecolor{corString}{RGB}{163,21,21}
\definecolor{corErro}{RGB}{180,0,0}
\definecolor{corOK}{RGB}{0,120,0}

% ---------- listings: mapeamento de acentos/símbolos (literate) ----------
\lstset{
  basicstyle=\ttfamily\footnotesize,
  backgroundcolor=\color{corFundo},
  frame=single,
  rulecolor=\color{corMoldura},
  breaklines=true,
  breakatwhitespace=false,
  columns=fullflexible,
  keepspaces=true,
  showstringspaces=false,
  upquote=true,
  extendedchars=true,
  literate=%
    {á}{{\'a}}1 {à}{{\`a}}1 {â}{{\^a}}1 {ã}{{\~a}}1 {ä}{{\"a}}1
    {é}{{\'e}}1 {ê}{{\^e}}1 {è}{{\`e}}1 {ë}{{\"e}}1
    {í}{{\'i}}1 {î}{{\^i}}1 {ï}{{\"i}}1
    {ó}{{\'o}}1 {ô}{{\^o}}1 {õ}{{\~o}}1 {ö}{{\"o}}1
    {ú}{{\'u}}1 {û}{{\^u}}1 {ü}{{\"u}}1
    {ç}{{\c c}}1 {ñ}{{\~n}}1
    {Á}{{\'A}}1 {À}{{\`A}}1 {Â}{{\^A}}1 {Ã}{{\~A}}1
    {É}{{\'E}}1 {Ê}{{\^E}}1 {È}{{\`E}}1
    {Í}{{\'I}}1 {Ó}{{\'O}}1 {Ô}{{\^O}}1 {Õ}{{\~O}}1
    {Ú}{{\'U}}1 {Ç}{{\c C}}1 {Ñ}{{\~N}}1
    {º}{{\textordmasculine}}1 {ª}{{\textordfeminine}}1 {°}{{\textdegree}}1
    {§}{{\S}}1
    {–}{{--}}1 {—}{{---}}2
    {“}{{``}}1 {”}{{''}}1 {‘}{{`}}1 {’}{{'}}1
    {…}{{\ldots}}1 {×}{{$\times$}}1
}

\lstdefinestyle{terminal}{}
\lstdefinestyle{json}{}
\lstdefinestyle{py}{
  language=Python,
  keywordstyle=\color{corChave}\bfseries,
  commentstyle=\color{corComent}\itshape,
  stringstyle=\color{corString},
  morekeywords={dataclass,field},
}
\lstdefinestyle{html}{
  language=HTML,
  keywordstyle=\color{corChave},
  stringstyle=\color{corString},
  commentstyle=\color{corComent}\itshape,
}

\newcommand{\status}[1]{\texttt{#1}}

% =====================================================================
\title{\textbf{Verificador de Contestação --- Fase 1}\\[2mm]
\large Relatório de execução: arquitetura, funcionamento e saída}
\author{Camada determinística (sem LLM)\\
\small Repositório: \url{https://github.com/gbsgubi/verificador-contestacao}}
\date{\today}

\begin{document}
\maketitle

\begin{abstract}
\noindent Este documento descreve o script de verificação de \textbf{minutas de
contestação} em ações previdenciárias de aposentadoria especial, construído na
\emph{Fase 1} (camada determinística, sem qualquer modelo de linguagem). São
apresentados: a arquitetura e o funcionamento do script, a forma de execução,
os \textbf{exemplos de minuta} usados como entrada (HTML nativo do SAPIENS e o
texto extraído do PDF) e, sobretudo, a \textbf{saída} produzida --- nas
representações de terminal e JSON --- com a interpretação do veredito.
\end{abstract}

\tableofcontents
\bigskip

% =====================================================================
\section{Visão geral}

A ferramenta recebe uma minuta de contestação (gerada no SAPIENS a partir do
modelo padrão n\textsuperscript{o} 544316, nativamente HTML e tipicamente
exportada em PDF) e devolve um \textbf{diagnóstico de aprovação ou reprovação},
apontando os problemas, antes de o estagiário enviá-la ao procurador.

A arquitetura final prevê duas camadas: uma \textbf{determinística} (regras
explícitas, auditáveis e testáveis) e outra de \textbf{raciocínio} (LLM). Este
relatório trata exclusivamente da camada determinística, que cobre verificações
de \emph{presença/ausência} e de \emph{padrão textual}.

\subsection{O que esta fase verifica}
\begin{itemize}[leftmargin=1.4em]
  \item \textbf{Bloco 1 --- Estrutural:} tipo de endereçamento (JEF
        $\times$ Justiça Federal comum); campos obrigatórios (NB, DER, número do
        processo, nome do autor, tipo de benefício); marcações de edição não
        removidas (colchetes, \emph{placeholders}, instruções imperativas e
        destaque amarelo no HTML).
  \item \textbf{Bloco 2 --- Preliminares (presença):} detecta se cada uma das
        dez preliminares está \status{PRESENTE} ou \status{AUSENTE}; e a exceção
        determinística JEF + Renúncia aos 60 SM ausente $\Rightarrow$
        \status{ERRO}.
  \item \textbf{Bloco 6 --- Redistribuição:} sinaliza (nunca como erro) os
        perfis \status{VIGILANTE}, \status{SAÚDE} e \status{PROFESSOR}.
  \item \textbf{Teses (Blocos 3/4/5):} apenas cataloga as teses padronizadas
        presentes (\texttt{CTN-}, \texttt{DIVESP-NOTA-}, \ldots).
\end{itemize}

% =====================================================================
\section{Arquitetura e funcionamento}

O processamento é um \emph{pipeline} de quatro etapas, orquestrado por
\texttt{src/verificador.py}:

\begin{enumerate}[leftmargin=1.6em]
  \item \textbf{Extração} (\texttt{extractor.py}): \texttt{pdfplumber} para PDF e
        \texttt{BeautifulSoup} para HTML. No HTML, os destaques amarelos (classe
        CSS \texttt{amarelo} ou \texttt{style} com fundo amarelo) são detectados
        diretamente; no PDF, o amarelo vira texto comum e a detecção é apenas
        textual.
  \item \textbf{Segmentação} (\texttt{segmenter.py}): divide o texto em seções
        (\textsc{síntese}, \textsc{preliminares}, \textsc{mérito}, \ldots) e
        extrai os blocos de período/agente, produzindo o objeto \texttt{Minuta}.
  \item \textbf{Regras} (\texttt{src/rules/}): cada verificação implementa a
        interface comum \texttt{Regra} e é coletada por um \emph{registry}.
  \item \textbf{Relatório} (\texttt{report.py}): agrega os \texttt{Resultado}s,
        calcula o veredito e serializa em JSON e em texto para o terminal.
\end{enumerate}

\begin{table}[h]
\centering
\small
\begin{tabular}{@{}lp{10.2cm}@{}}
\toprule
\textbf{Arquivo} & \textbf{Responsabilidade} \\
\midrule
\texttt{src/utils.py}            & Normalização de texto (sem acentos, minúsculas) e busca. \\
\texttt{src/extractor.py}        & Extração de PDF/HTML; detecção de amarelo no HTML. \\
\texttt{src/segmenter.py}        & Segmentação em seções; objeto \texttt{Minuta}. \\
\texttt{src/report.py}           & \texttt{Resultado}/\texttt{Relatorio}; veredito; JSON e terminal. \\
\texttt{src/verificador.py}      & Orquestrador do pipeline + interface de linha de comando. \\
\texttt{src/rules/\_\_init\_\_.py} & Interface \texttt{Regra} + \emph{registry} de regras. \\
\texttt{src/rules/estrutural.py} & Bloco 1. \\
\texttt{src/rules/preliminares.py} & Bloco 2 (presença + exceção JEF). \\
\texttt{src/rules/redistribuicao.py} & Bloco 6. \\
\texttt{src/rules/teses.py}      & Catalogação de teses. \\
\texttt{src/rules/llm\_stubs.py} & Pontos de extensão para a fase do LLM. \\
\bottomrule
\end{tabular}
\caption{Módulos do projeto e suas responsabilidades.}
\end{table}

\subsection{A interface comum de regra}
Toda verificação --- determinística agora ou baseada em LLM no futuro --- segue
a mesma interface, o que permite plugar novas regras sem alterar o orquestrador:

\begin{lstlisting}[style=py]
class Regra(ABC):
    id: str = ""          # identificador estável (ex.: "1.3-marcacoes-edicao")
    bloco: str = ""       # rótulo do bloco, para agrupar no relatório
    descricao: str = ""   # o que a regra checa
    requer_llm: bool = False

    @abstractmethod
    def verificar(self, minuta: Minuta) -> list[Resultado]:
        """Aplica a regra à minuta e devolve uma lista de resultados."""
        raise NotImplementedError


def registrar(cls):           # decorador que registra a regra no registry
    _REGISTRO.append(cls)
    return cls
\end{lstlisting}

% =====================================================================
\section{Como executar}

\subsection{Demonstração de um clique}
O script \texttt{demo.ps1} (Windows/PowerShell) cria o ambiente virtual se
necessário, instala as dependências e roda o verificador sobre a minuta de teste
(PDF e HTML):

\begin{lstlisting}[style=terminal]
powershell -ExecutionPolicy Bypass -File .\demo.ps1
\end{lstlisting}

\subsection{Linha de comando}
\begin{lstlisting}[style=terminal]
# relatório legível no terminal
python -m src.verificador tests/fixtures/minuta_teste.pdf

# formato HTML nativo do SAPIENS
python -m src.verificador tests/fixtures/minuta_teste.html

# gravar o relatório em JSON
python -m src.verificador tests/fixtures/minuta_teste.pdf --json saida.json
\end{lstlisting}

\noindent O \textbf{código de saída} é \texttt{1} quando o veredito é
\status{REPROVADO} e \texttt{0} quando \status{APROVADO} --- útil em automações.

% =====================================================================
\section{Exemplos de minuta (entrada)}

A minuta de teste é \emph{fictícia}: simula um caso no \textbf{JEF}, com dois
períodos especiais requeridos (ruído de 06/03/1997 a 31/12/2002 e calor de
01/01/2003 a 18/11/2010), e contém erros plantados de propósito --- em especial
uma anotação de edição em amarelo não removida.

\subsection{Formato HTML (nativo do SAPIENS)}
\begin{lstlisting}[style=html]
<!DOCTYPE html>
<html lang="pt-BR">
<head>
<meta charset="UTF-8">
<style>
  body { font-family: "Times New Roman", Times, serif; font-size: 12pt; }
  .titulo { text-align: center; font-weight: bold; text-transform: uppercase; }
  .secao { font-weight: bold; text-transform: uppercase; }
  .subsecao { font-weight: bold; }
  table { border-collapse: collapse; width: 100%; }
  td, th { border: 1px solid #000; padding: 5px 7px; vertical-align: top; }
  .rotulo { font-weight: bold; width: 22%; background: #f0f0f0; }
  .amarelo { background: #fff200; }   /* destaque a ser removido */
</style>
</head>
<body>

<p class="indent">EXCELENTÍSSIMO(A) SENHOR(A) DOUTOR(A) JUIZ(A) FEDERAL DA 10ª
VARA FEDERAL DO JUIZADO ESPECIAL FEDERAL DA SUBSEÇÃO JUDICIÁRIA DE SÃO PAULO - SP</p>

<p class="right">Processo nº 5001234-56.2024.4.03.6183<br>
NUP: 01032.456789/2024-11<br>
NB: 46/123.456.789-0<br>
DER: 14/02/2022</p>

<p class="indent">O <strong>INSTITUTO NACIONAL DO SEGURO SOCIAL - INSS</strong>,
autarquia federal já qualificada nos autos do processo em epígrafe, que lhe move
<strong>JOÃO DA SILVA SANTOS</strong>, vem, respeitosamente, [...] apresentar</p>

<p class="titulo">CONTESTAÇÃO</p>

<p class="secao">I - SÍNTESE DA DEMANDA</p>
<p class="indent">Trata-se de ação em que pleiteia a parte autora a concessão do
benefício de aposentadoria especial após a conversão de tempo especial em comum,
alegando exposição a agentes nocivos nos períodos de 06/03/1997 a 31/12/2002 e de
01/01/2003 a 18/11/2010.</p>

<p class="secao">II - PRELIMINARES</p>

<p class="subsecao">II.1 - Do Juízo 100% Digital</p>
<p class="indent">O INSS concorda com a tramitação na modalidade "Juízo 100%
Digital" [...] (art. 183, §1º, do CPC c/c art. 5º, da Lei nº 11.419/2006).</p>

<p class="subsecao">II.2 - Da Audiência de Conciliação</p>
<p class="indent">[...] o INSS informa não possuir interesse na realização de
audiência de conciliação (art. 334, §4º, II, CPC/2015).</p>

<p class="subsecao">II.3 - Da Renúncia ao Valor Excedente a 60 Salários-Mínimos</p>
<p class="indent">Tratando-se de feito que tramita perante o Juizado Especial
Federal, é necessário que a parte autora renuncie expressamente ao montante que
ultrapasse sessenta salários mínimos (art. 3º, §2º, da Lei n.º 10.259/2001) [...]</p>

<p class="subsecao">II.4 - Da Decadência</p>
<p class="indent">Considerando o fato de ter decorrido mais de 10 (dez) anos [...]
decaiu o direito da parte autora de revisar o ato de concessão de seu benefício.</p>

<!-- ERRO PLANTADO: destaque amarelo (instrução de edição) não removido -->
<p class="amarelo">[VERIFICAR: deixar a preliminar de decadência somente quando
contabilizar 10 anos da data de ajuizamento até a presente data]</p>

<p class="secao">III - DO MÉRITO</p>

<p class="subsecao">III.1 - Período de 06/03/1997 a 31/12/2002 - Agente: Ruído</p>
<table>
  <tr><td class="rotulo">Período(s):</td><td>06/03/1997 a 31/12/2002</td></tr>
  <tr><td class="rotulo">Prova apresentada:</td>
      <td>PPP de fls. 45/46 do PA (data de emissão: 12/03/2019)</td></tr>
  <tr><td class="rotulo">Fundamentos da defesa:</td><td>
    <strong>Questões Prejudiciais / Vícios Formais do PPP:</strong><br>
    1. A parte autora não comprova que o vistor do PPP possua poderes de
       representação da empresa.<br>
    <strong>Agentes Nocivos:</strong><br>
    1. Exposição ao ruído dentro do limite de tolerância - 90 dB(A).<br>
    2. Não há responsável técnico pelos registros ambientais no período [...]<br>
    3. A metodologia de avaliação informada não atende à legislação. Para
       períodos entre 03/12/1998 e 18/11/2003, [...] devem obediência à NR-15.
  </td></tr>
</table>

<p class="subsecao">III.2 - Período de 01/01/2003 a 18/11/2010 - Agente: Calor</p>
<table>
  <tr><td class="rotulo">Período(s):</td><td>01/01/2003 a 18/11/2010</td></tr>
  <tr><td class="rotulo">Prova apresentada:</td>
      <td>PPP de fls. 45/46 do PA (data de emissão: 12/03/2019)</td></tr>
  <tr><td class="rotulo">Fundamentos da defesa:</td><td>
    <strong>Agentes Nocivos:</strong><br>
    1. A exposição é inferior a 25ºC (patamar mínimo do Anexo 3 da NR-15 [...]).<br>
    2. Não há responsável técnico [...] exigível a partir de 14/10/1996.<br>
    3. As medições devem ser feitas em "IBUTG" (Anexo 3 da NR-15).
  </td></tr>
</table>

<p class="secao">IV - OUTROS FUNDAMENTOS</p>
<p class="indent">Fica vedada a conversão de tempo especial em comum para períodos
posteriores a 13/11/2019 (Art. 25, §2º, da EC nº 103/2019 [...]).</p>

<p class="secao">V - DOS PEDIDOS</p>
<p>a) o acolhimento das preliminares arguidas;</p>
<p>b) no mérito, a total improcedência dos pedidos; [...]</p>

<p class="assinatura">_______________________________________<br>
Procurador Federal</p>
</body>
</html>
\end{lstlisting}

\noindent\textit{Observação:} o trecho com a classe \texttt{amarelo} é o erro
plantado de \emph{alta prioridade} --- uma instrução de edição que deveria ter
sido removida antes do envio.

\subsection{Texto extraído do PDF}
A mesma minuta em PDF, após a extração com \texttt{pdfplumber}, produz o texto
plano abaixo (o destaque amarelo deixa de existir, mas o padrão textual
--- colchetes e a instrução ``VERIFICAR'' / ``deixar a preliminar'' --- permanece
detectável):

\begin{lstlisting}[style=terminal]
EXCELENTÍSSIMO(A) SENHOR(A) DOUTOR(A) JUIZ(A) FEDERAL DA 10ª VARA FEDERAL DO
JUIZADO ESPECIAL FEDERAL DA SUBSEÇÃO JUDICIÁRIA DE SÃO PAULO - SP
Processo nº 5001234-56.2024.4.03.6183
NUP: 01032.456789/2024-11
NB: 46/123.456.789-0
DER: 14/02/2022
O INSTITUTO NACIONAL DO SEGURO SOCIAL - INSS [...], que lhe move JOÃO DA SILVA
SANTOS, [...] apresentar CONTESTAÇÃO [...]

I - SÍNTESE DA DEMANDA
Trata-se de ação em que pleiteia a parte autora a concessão do benefício de
aposentadoria especial após a conversão de tempo especial em comum, alegando
exposição a agentes nocivos nos períodos de 06/03/1997 a 31/12/2002 e de
01/01/2003 a 18/11/2010.

II - PRELIMINARES
II.1 - Do Juízo 100% Digital [...]
II.2 - Da Audiência de Conciliação [...]
II.3 - Da Renúncia ao Valor Excedente a 60 Salários-Mínimos [...]
II.4 - Da Decadência [...] decaiu o direito da parte autora de revisar o ato de
concessão de seu benefício.
[VERIFICAR: deixar a preliminar de decadência somente quando contabilizar 10
anos da data de ajuizamento até a presente data]

III - DO MÉRITO
III.1 - Período de 06/03/1997 a 31/12/2002 - Agente: Ruído
Período(s): 06/03/1997 a 31/12/2002
Prova apresentada: PPP de fls. 45/46 do PA (data de emissão: 12/03/2019)
Fundamentos da defesa: [...] 1. Exposição ao ruído dentro do limite de
tolerância - 90 dB(A). [...]
III.2 - Período de 01/01/2003 a 18/11/2010 - Agente: Calor [...]

IV - OUTROS FUNDAMENTOS [...]
V - DOS PEDIDOS [...]
São Paulo, 10 de março de 2025.
_______________________________________
Procurador Federal
\end{lstlisting}

% =====================================================================
\section{Saída da execução}

Esta é a parte central. A ferramenta produz duas representações do mesmo
relatório: uma \textbf{legível no terminal} (agrupada por bloco, com veredito no
topo) e uma \textbf{JSON} estruturada.

\subsection{Representação no terminal}
\begin{lstlisting}[style=terminal]
======================================================================
  VEREDITO: REPROVADO
  Erros: 1   |   Pontos de verificação manual: 0
======================================================================

### Bloco 1 - Estrutural
  [INFO] Tipo de endereçamento (JEF x Justiça Federal comum)
      Endereçamento identificado: JEF (Juizado Especial Federal).
      evidência: ...feito que tramita perante o Juizado Especial Federal...
  [OK] Campo obrigatório: NB
      NB presente.
      evidência: ...NUP: 01032.456789/2024-11 NB: 46/123.456.789-0 DER: 14/02/2022...
  [OK] Campo obrigatório: DER
      DER presente.
  [OK] Campo obrigatório: Número do processo
      Número do processo presente.
      evidência: ...Processo nº 5001234-56.2024.4.03.6183 NUP: 01032.456789/2024-11...
  [OK] Campo obrigatório: Nome do autor
      Nome do autor presente.
      evidência: ...que lhe move JOÃO DA SILVA SANTOS, vem, respeitosamente...
  [OK] Campo obrigatório: Tipo de benefício
      Tipo de benefício presente.
      evidência: ...a concessão do benefício de aposentadoria especial...
  [ERRO] Marcações de edição não removidas (colchetes, placeholders, instruções, amarelo)
      Há 2 marcação(ões) de edição não removida(s) na minuta — devem ser
      eliminadas antes do envio.
      evidência: instrução remanescente: "...[VERIFICAR: deixar a preliminar de
      decadência some..." | instrução remanescente: "...deixar a preliminar de
      decadência somente quando contabiliz..."

### Bloco 2 - Preliminares (presença)
  [INFO] Preliminar: Juízo 100% Digital            PRESENTE.
  [INFO] Preliminar: Audiência de Conciliação      PRESENTE.
  [INFO] Preliminar: Renúncia aos 60 Salários-Mínimos   PRESENTE.
  [INFO] Preliminar: Decadência                    PRESENTE.
  [INFO] Preliminar: Prescrição                    AUSENTE.
  [INFO] Preliminar: Coisa Julgada                 AUSENTE.
  [INFO] Preliminar: Litispendência                AUSENTE.
  [INFO] Preliminar: PPP não apresentado administrativamente   AUSENTE.
  [INFO] Preliminar: Períodos reconhecidos administrativamente AUSENTE.
  [INFO] Preliminar: Petição inicial inepta        AUSENTE.
  [OK] Renúncia aos 60 SM obrigatória quando o juízo é o JEF
      Processo no JEF e preliminar de Renúncia aos 60 SM presente.

### Bloco 6 - Redistribuição
  [OK] Indício de categoria vigilante/vigia/guarda/policial
      Sem indício de categoria vigilante/vigia/guarda/policial.
  [OK] Indício de profissional/ambiente de saúde + agente biológico
      Sem indício combinado de profissional de saúde e agente biológico.
  [OK] Indício de professor com período posterior a 1981
      Sem indício de professor em período posterior a 1981.

### Teses presentes (Blocos 3/4/5 - catalogação)
  [INFO] Catalogação das teses padronizadas presentes (CTN-/DIVESP- etc.)
      Nenhuma tese padronizada (CTN-/DIVESP-) encontrada no texto. A minuta de
      teste usa fundamentos em prosa, sem códigos de tese.

----------------------------------------------------------------------
  VERIFICAÇÕES ADIADAS PARA A FASE DO LLM (não executadas):
   - [3.x-compat-tese-agente-periodo] Compatibilidade tese × agente × período
   - [2.8-situacoes-ppp] Distinção entre as três situações de PPP não apresentado
   - [2.45-decadencia-prescricao-datas] Cabimento de decadência/prescrição por datas
   - [1.4-cobertura-periodos-inicial] Cobertura dos períodos requeridos na inicial
\end{lstlisting}

\subsection{Representação JSON}
Cada verificação é um objeto
\texttt{\{ id, bloco, descricao, status, mensagem, evidencia \}}, com
\texttt{status} $\in$ \{\status{OK}, \status{ERRO}, \status{VERIFICAR},
\status{INFO}\}. Abaixo, o cabeçalho e os itens mais relevantes (saída
abreviada):

\begin{lstlisting}[style=json]
{
  "veredito": "REPROVADO",
  "total_erros": 1,
  "total_verificar": 0,
  "verificacoes": [
    {
      "id": "1.1-enderecamento",
      "bloco": "Bloco 1 - Estrutural",
      "descricao": "Tipo de endereçamento (JEF x Justiça Federal comum)",
      "status": "INFO",
      "mensagem": "Endereçamento identificado: JEF (Juizado Especial Federal).",
      "evidencia": "...tramita perante o Juizado Especial Federal, é necessário..."
    },
    {
      "id": "1.3-marcacoes-edicao",
      "bloco": "Bloco 1 - Estrutural",
      "descricao": "Marcações de edição não removidas (colchetes, placeholders, instruções, amarelo)",
      "status": "ERRO",
      "mensagem": "Há 2 marcação(ões) de edição não removida(s) na minuta — devem ser eliminadas antes do envio.",
      "evidencia": "instrução remanescente: \"...[VERIFICAR: deixar a preliminar de decadência some...\" | ..."
    },
    {
      "id": "2.presenca:prescricao",
      "bloco": "Bloco 2 - Preliminares (presença)",
      "descricao": "Preliminar: Prescrição",
      "status": "INFO",
      "mensagem": "AUSENTE.",
      "evidencia": null
    },
    {
      "id": "2.3-renuncia-jef",
      "bloco": "Bloco 2 - Preliminares (presença)",
      "descricao": "Renúncia aos 60 SM obrigatória quando o juízo é o JEF",
      "status": "OK",
      "mensagem": "Processo no JEF e preliminar de Renúncia aos 60 SM presente.",
      "evidencia": null
    }
  ]
}
\end{lstlisting}

% =====================================================================
\section{Interpretação dos resultados}

A regra do veredito é simples: \textbf{\status{REPROVADO}} se houver qualquer
\status{ERRO}; caso contrário \textbf{\status{APROVADO}}, listando os
\status{VERIFICAR} pendentes. Na minuta de teste o veredito é
\status{REPROVADO} por causa de \emph{um} erro: a marcação de edição não
removida.

\begin{table}[h]
\centering
\small
\begin{tabular}{@{}llp{7.6cm}@{}}
\toprule
\textbf{Item} & \textbf{Status} & \textbf{Leitura} \\
\midrule
Endereçamento & \status{INFO}  & Identificado como JEF, coerente com o nº do processo. \\
Campos obrigatórios & \status{OK} & NB, DER, processo, autor e benefício presentes. \\
Marcação de edição & \status{ERRO} & Instrução ``VERIFICAR/decadência'' não removida. \\
Preliminares & \status{INFO} & 4 presentes; Prescrição e outras ausentes (fato, não juízo). \\
Renúncia 60 SM (JEF) & \status{OK} & Presente, como exigido no JEF. \\
Redistribuição & \status{OK} & Sem perfis VIGILANTE/SAÚDE/PROFESSOR. \\
Teses padronizadas & \status{INFO} & Nenhuma (a minuta usa prosa). \\
\bottomrule
\end{tabular}
\caption{Resumo da leitura do relatório sobre a minuta de teste.}
\end{table}

\subsection{O que ficou para a fase do LLM}
Algumas verificações exigem raciocínio contextual e foram deixadas como
\emph{stubs} documentados (mesma interface \texttt{Regra}), portanto
\textbf{não} executadas nesta fase:
\begin{itemize}[leftmargin=1.4em]
  \item compatibilidade entre tese, período e agente nocivo (marcos temporais
        de ruído/calor) --- p.\,ex.\ a tese de ruído na fronteira de 03/12/1998;
  \item distinção entre as três situações de PPP não apresentado;
  \item cabimento de decadência/prescrição segundo datas e tipo de ação
        (concessão $\times$ revisão);
  \item cobertura dos períodos requeridos na petição inicial.
\end{itemize}

Por isso, propositalmente, esta fase \emph{não} sinaliza a incompatibilidade da
tese de ruído nem a possível decadência indevida --- ambas dependem da camada de
LLM.

\end{document}
